
%The error-aware structure of revision traces---supervising the trajectory from mistake to recovery, not just the correct endpoint---appears to be the key ingredient: SRT's largest gains are on out-of-distribution benchmarks where error recovery is most valuable, and SFT on expert traces actually degrades performance where the distribution gap is largest. Self-distillation then compresses this two-attempt behavior into shorter, single-attempt reasoning while preserving accuracy and revision capability (Figure~\ref{fig:self-revision}). Whether the resulting model genuinely anticipates errors or merely produces shorter outputs remains an open empirical question.


\section{Discussion}
\label{sec:discussion}
% \looseness-1\ours{} shows that a model's own revision capability, once elicited with a small amount of supervised training on self-generated revisions, can transform sparse binary outcome feedback into dense token-level supervision without requiring an external teacher. Across two model families and eight benchmarks, this yields consistent gains of at least \(10\%\) over the base models and outperforms all baselines under the same data budget (\Cref{tab:main_results}); independent experiments on a different task, model, and pipeline variant corroborate both the findings and key design principles (\Cref{app:countdown}).
% More broadly, our results suggest a different view of self-improvement: rather than relying on stronger teacher-generated reasoning responses, a model can learn to convert its own reward-conditioned revisions into token-level supervision.


\looseness-1 \textbf{Conclusion.} \ours{} shows that a model’s revision ability, elicited with a small amount of supervised training on self-generated revisions, can transform sparse binary outcome reward into dense token-level supervision without requiring an external teacher or high-quality demonstrations. 
Across two model families and eight benchmarks, this yields consistent \(10\%+\) gains over the base models and outperforms all baselines under the same training sample budget (\Cref{tab:main_results}). 
We also show experiments on more broader tasks and models in Appendix \ref{app:countdown}. 
Extensive ablations and analyses show that \ours{} also leads to improved revision capability which can lead to further gains on accuracy.
We hope that by eliminating the need for high-quality demonstrations, we enable future work to apply self-distillation to broader domains.

% \textbf{Concurrent and Related Works.} We discuss related work in detail in \Cref{sec:related}; here we highlight the closest concurrent directions. On-policy Self-Distillation (OPSD) \citep{zhao2026selfdistilledreasoneronpolicyselfdistillation} and Self-Distillation Fine-Tuning (SDFT) \citep{shenfeld2026self} are closely related in that they use the model itself to provide dense token-level supervision. However, both assume access to stronger teacher-generated solutions or gold solutions. By contrast, \ours{} requires neither an external teacher nor gold solutions, and instead derives token-level supervision directly from sparse outcome reward. Empirically, we find that SDFT can substantially underperform \ours{} despite its stronger supervision, suggesting that using gold solutions may be brittle when the reasoning style of those solutions differs from the model's own responses.

%\looseness-1Self-Distillation Policy Optimization (SDPO) \citep{hubotter2026reinforcement} also avoids an external teacher, but assumes access to finer-grained signals indicating where the initial response fails. Such signals may come from successful self-generated responses filtered by reward, which can be expensive to collect and reuse as supervision, or from environment feedback available only in highly structured settings such as coding. 
%In contrast, \ours{} requires only binary outcome reward, whether the initial response is correct or incorrect, making it substantially more lightweight and broadly applicable.
% \looseness-1Self-Distillation Policy Optimization (SDPO) \citep{hubotter2026reinforcement}  also avoids an external teacher, but assumes access to finer-grained signals such as environment feedback (e.g., runtime errors or test results) or successful self-generated solutions filtered by reward, which can be expensive to collect. In contrast, \ours{}'s reviser conditions only on the student's original attempt and its binary outcome reward, requiring neither rich environment feedback nor successful reference solutions.


\looseness-1\textbf{Limitations and Open Directions.} Our study focuses on instruct models that generate  short and concise responses, consistent with prior self-distillation works. An important next step is extending self-distillation to \emph{thinking models}, which produce long, exploratory chains of thought. This is challenging because such responses may include false starts and partial corrections that are not mistakes, making it difficult to distinguish productive exploration from genuine errors and assign credit beyond local token decisions. We provide preliminary evidence of this challenge in \Cref{app:opsd-thinking}: applying SDFT to Qwen3-4B with thinking enabled during training hurts performance on competition math benchmarks.

\looseness-1Furthermore, we focus on verifiable domains such as math and coding. Extending \ours{} to domains without verifiable rewards remains an open problem. One promising direction is to define rewards using meta-cognitive signals \citep{didolkar2024metacognitive,didolkar2025metacognitive,shao2025dr}, such as consistency or self-correction. We leave this to future work.



\iffalse
\textbf{Limitations and Open Directions.} Our study focuses on instruct models that produce relatively short, concise reasoning before the final answer, which is also the regime considered by existing self-distillation methods. A natural and important next step is to extend self-distillation to \emph{thinking models}, where the model generates long, exploratory chains of thought before reaching an answer. This is not a straightforward extension: unlike short-form reasoning, long chains of thought may contain false starts, partial corrections, and exploratory branches that are not themselves errors but part of successful reasoning. Developing supervision strategies that can distinguish productive exploration from genuine mistakes, and assign credit at the level of the overall reasoning trajectory rather than local token decisions, remains an important open problem.
We provide preliminary evidence for this challenge in \Cref{app:opsd-thinking}, where applying SDFT to Qwen3-4B with thinking enabled during training degrades performance on competition math benchmarks.

Furthermore, we focus on domains such as math and coding, where reward is based on correctness of the final answer. Extending self-distillation framework to domains without verifiable outcomes remains another important open question. One promising direction is to define rewards in terms of meta-cognitive skills \citep{didolkar2024metacognitive,didolkar2025metacognitive,shao2025dr}; for example, consistency, self-correction, or other indicators of reasoning quality. We leave this direction to future work.
\fi 


% \textbf{Concurrent and Related Works:} We provide a detailed discussion of related work in \Cref{sec:related}, focusing on on-policy distillation and self-improvement methods. Here, we briefly compare our method to some of the most recent concurrent works.

% On-policy Self-Distillation (OPSD) \citep{zhao2026selfdistilledreasoneronpolicyselfdistillation} and Self-Distillation Fine-Tuning (SDFT) \citep{shenfeld2026self} the same model under training to provide dense token-level supervision to itself. However, both rely on access to gold solutions, possibly generated from a stronger teacher. In contrast, \ours{} requires no external teacher or gold solutions, and instead derives token-level supervision directly from sparse outcome reward. Empirically, we find that OPSD can substantially underperform \ours{} despite having access to gold solutions, suggesting that mismatches between reasoning style of gold solutions and the model's own responses can make self-distillation ineffective.

% \looseness-1Self-Distillation Policy Optimization (SDPO) \citep{hubotter2026reinforcement} also avoids an external teacher, but assumes access to finer-grained error signals indicating where the initial response fails. Such signals may come from successful self-generated responses filtered by reward, which can be expensive to collect and use as feedback, or from environment feedback available only in very structured domains such as coding, which can be difficult to create in general domains. By comparison, \ours{} requires only binary outcome supervision, whether the initial response is correct or incorrect, and still converts this signal into dense token-level guidance, making it substantially more lightweight and broadly applicable.


% \textbf{Limitations and Open Directions.} Our study focuses on instruct models that produce relatively short, concise reasoning before the final answer, which is also the regime considered by existing self-distillation methods. A natural and important next step is to extend self-distillation to \emph{thinking models}, where the model generates long, exploratory chains of thought before reaching an answer. This is not a straightforward extension: unlike short-form reasoning, long chains of thought may contain multiple false starts, partial corrections, and exploratory branches that are not themselves errors but part of successful reasoning. Developing supervision strategies that can distinguish productive exploration from mistakes and assign credit at the level of the overall reasoning remains an important open problem.

% Furthermore, we focus on domains such as math and coding, where correctness of the final answer provides a natural reward signal. Extending our approach to domains without such verifiable rewards remains an important open question. One promising direction is to define rewards in terms of meta-cognitive skills; for example, consistency, self-correction, or other indicators of reasoning quality, and we leave this exploration to future work.


% \textbf{Limitations and Open directions.} In this work, we primarily focus on instruct model, models that generate a short  and concise reasoning before generating the final answer and have been utilized by all recent works on self-distillation. One key direction is to extend self-distillation to thinking models, where the model produces extended and long chains of thought before generating the final attempt. Critical challenges include how should we come up with the right strategy to provide supervision with long chain of thought steps, which are extremely exploratory in structure and can contain several incorrect attempts, in such a way that they affect the final reasoning of the model towards the final answer. 



% Extending \ours{} to thinking models---where the model produces extended, structured chains of thought---is nontrivial: per-token KL divergence requires teacher and student distributions to be comparable at each position, and in thinking models the teacher's distribution conditioned on $(Q, A, P_r)$ can diverge sharply from the student's exploratory reasoning. 
% This difficulty extends to other on-policy distillation methods: OPSD does not work out of the box for thinking models, and even vanilla OPD, while applicable, can be sensitive to model capacity and training configuration. Adapting token-level distillation to structured thinking formats is an important direction we plan to pursue.
% A separate limitation of \ours{} is that our method also requires the base model to self-correct at a nontrivial rate---it amplifies existing revision capability rather than creating it. And Phase~2 gains are uneven across domains: strong for code ($+8.2$ on LiveCodeBench for \qwen{}) but zero or negative on some math benchmarks, for reasons we do not yet fully understand. \sk{todo: can add exp result (maybe in appendix or something) where opsd fails for thinking models}

%This difficulty is shared by on-policy distillation methods in general\edward{Is this true? IIRC and in my (limited) experience, OPD works reasonably well for thinking models}\sk{i can reword it more carefully, was thinking about the 8b --> 4b case hurting 4b's performance sometimes; i guess the high level point i'm trying to make is that further training ``thinking'' models that are already crammed is hard / the models are very sensitive}, and adapting token-level distillation to structured thinking formats is an important direction we plan to pursue.

% \textbf{Open questions.} Several directions remain open: whether iterated teacher synchronization sustains gains beyond the one additional round we tested (\Cref{fig:evolution}), whether richer reward signals yield better revision teachers than binary reward, whether \ours{}'s advantage persists at larger model scale, and how to adapt on-policy (self) distillation to thinking models where token-level KL is no longer a natural objective. \sk{will reword last sentence; not sure if ``natural objective'' is the right phrasing here}
